\documentclass[11pt,a4paper,twocolumn]{article}
\usepackage[utf8]{inputenc}
\usepackage[slovak]{babel}
\usepackage[T1]{fontenc}
\usepackage[margin=2cm]{geometry}
\usepackage{graphicx}
\usepackage{cite}
\usepackage{hyperref}
\usepackage{booktabs}

\title{\textbf{Pasívny audit sieťovej bezpečnosti a AI filtrácia nepresností v kontexte smernice NIS2}}

\author{
    \textbf{Lukáš Štefan, Jakub Šima, Vincent Gavenda,}\\
    \textbf{Rony Šoóš, Oliver Šutara}\\
    \vspace{0.3em}\\
    Fakulta informatiky a informačných technológií STU v Bratislave\\
    \texttt{xstefan@stuba.sk}
}

\date{\today}

\begin{document}

\maketitle

\begin{abstract}
Európska smernica NIS2 zavádza prísne požiadavky na kybernetickú bezpečnosť pre tisíce stredne veľkých podnikov. Tradičné metódy auditu však trpia vysokou chybovosťou alebo právnymi rizikami spojenými s invazívnymi skenmi. Tento prehľadový článok analyzuje súčasný stav v oblasti neinvazívneho auditu sieťových protokolov (DNS, DMARC, SSL/TLS) a navrhuje architektúru systému LANCE. Riešenie kombinuje pasívny zber verejných dát s AI triážou na zníženie miery falošných poplachov pod 10\,\% pri dodržaní pravidiel GDPR.
\end{abstract}

\section{Úvod}
Smernica NIS2 \cite{nis2directive} rozširuje povinnosti kybernetickej bezpečnosti na stredné podniky v kľúčových sektoroch. Organizácie čelia výzve, ako efektívne vyhodnotiť stav svojej infraštruktúry.

Cieľom tohto článku je zmapovať metódy pasívneho prieskumu a navrhnúť mechanizmus filtrácie nepresností pomocou umelej inteligencie.

\section{Súčasný stav v oblasti}
Pasívny prieskum spolieha na analýzu dát, ktoré cieľový systém verejne publikuje.

\subsection{Protokoly e-mailovej bezpečnosti}
Nastavenia DMARC, SPF a DKIM sú kľúčové pre prevenciu spoofingu. Podľa štúdií \cite{valimail2023, cloudflare2024} väčšina firiem nemá DMARC nastavený v režime presadzovania (\texttt{p=reject}).

\subsection{Bezpečnostné hlavičky a SSL/TLS}
Absencia HTTP hlavičiek ako HSTS či CSP \cite{mozillaobs, securityheaders} indikuje medzery v správe webových aplikácií.

\subsection{Právny rámec (GDPR)}
Pasívny zber DNS a verejných údajov z webov spĺňa podmienky článku 6 ods. 1 písm. f) GDPR \cite{gdpr2016} a nevyžaduje invazívne zásahy do infraštruktúry.

\section{Architektúra a metodika LANCE}
Navrhnutá architektúra pozostáva z modulu pasívneho skenovania, výpočtu skóre dôveryhodnosti a AI filtrácie.

\section{Vyhodnotenie a diskusia}
Porovnanie s existujúcimi nástrojmi ukazuje, že využitie LLM triáže výrazne znižuje prácnosť audítorov pri overovaní nálezov.

\section{Záver}
Navrhnutý prístup umožňuje bezpečné a presné vyhodnotenie sieťovej bezpečnosti podľa NIS2 bez nutnosti invazívnych testov.

\bibliographystyle{IEEEtran}
\bibliography{literatura}

\end{document}